% Generated from locale/tr/content/lambda-calculus/introduction/primitive-recursion.tex; do not hand-edit.


\olfileid[tr]{lam}{int}{pr}
\olsection{\usetoken{S}{lambda definable} Fonksiyonlar İlkel Özyineleme Altında Kapalıdır}

İlkel özyinelemeye gelince artık biraz çalışmamız gerekir. Aşamalar hâlinde
ilerleyeceğiz. Önceki gibi, sırasıyla $g$ ve $h$ fonksiyonlarını
!!{lambda define} eden $G$ ve $H$ terimlerine zaten sahip olduğumuzu
varsayarak, şu denklemlerle tanımlanan $f$ fonksiyonunu !!{lambda define}s
bir terim isteriz:
\begin{align*}
f(0, \vec z) & = g(\vec z) \\
f(x+1, \vec z) & = h(z, f(x,\vec z), \vec z).
\end{align*}
Dolayısıyla genel olarak $G'$ ve $H'$ lambda terimleri verildiğinde, her $n$
doğal sayısı için
\begin{align*}
F(\num{0}, \vec z) & \equiv G(\vec z) \\
F(\overline{n+1}, \vec z) & \equiv H(\num{n}, F(\num{n}, \vec z), \vec z)
\end{align*}
olacak biçimde bir $F$ terimi bulmak yeterlidir. $G'$ ile $H'$'nin $g$ ile
$h$'yi !!{lambda define} etmesi, $\vec z$ yerine $\num{\vec m}$ sayıları
koyduğumuzda $F(\num{n+1},\num{\vec m})$'nin doğru yanıta normalleşeceği
anlamına gelir.

Bunun içinse
\begin{align*}
F(\num 0) & \equiv G \\
F(\overline {n+1}) & \equiv H(\num{n},F(\num{n}))
\intertext{koşullarını her $n$ doğal sayısı için sağlayan bir $F$ bulmak yeterlidir; burada}
G & = \lambd[\vec z][G'(\vec z)] \text{ ve}\\
H(u,v) & = \lambd[\vec z][H'(u,v(u,\vec z),\vec z)].
\end{align*}
Başka bir deyişle lambda oyunuyla ek $\vec z$ parametrelerini dert etmekten
kaçınabiliriz; bunlar lambda gösteriminin içinde soğurulur.

$F$ terimini tanımlamadan önce sıralı ikilileri işleyecek bir düzeneğe
gereksinimimiz vardır. Bunu sonraki lemma sağlar.

\begin{lem}
Her $M$ ve $N$ lambda terimi çifti için
$D(M,N)(\num{0})\red M$ ve $D(M,N)(\num{1})\red N$ olacak biçimde bir $D$
lambda terimi vardır.
\end{lem}

\begin{proof}
Önce $K$ lambda terimini
\[
K(y) = \lambd[x][y]
\]
ile tanımlayın. Başka bir deyişle $K$, $\lambd[y][\lambd[x][y]]$ terimidir.
Başka açıdan bakarsak her $M$ için $K(M)$, her girdide $M$'yi döndüren sabit
bir fonksiyondur.

Şimdi $D(x,y,z)$'yi $D(x,y,z)=z(K(y))x$ ile tanımlayın. O zaman istendiği
gibi
\begin{align*}
D(M,N,\num 0) & \red \num 0 (K(N)) M \red M \text{ ve}\\
D(M,N,\num 1) & \red \num 1 (K(N)) M \red K(N) M \red N
\end{align*}
olur.
\end{proof}

Düşünce, $D(M,N)$'nin $\tuple{M,N}$ ikilisini temsil etmesi; $P$'nin böyle
bir ikiliyi temsil ettiği varsayılırsa $P(\num 0)$ ile $P(\num 1)$'in sol ve
sağ izdüşümleri $(P)_0$ ile $(P)_1$'i temsil etmesidir. Bundan sonra ikinci
gösterimleri kullanacağız.

\begin{lem}
!!{lambda definable} fonksiyonlar ilkel özyineleme altında kapalıdır.
\end{lem}

\begin{proof}
Herhangi $G$ ve $H$ terimleri verildiğinde, her $n$ doğal sayısı için
\begin{align*}
F(\num{0}) & \equiv G \\
F(\num{n+1}) & \equiv H(\num{n}, F(\num{n}))
\end{align*}
olacak biçimde bir $F$ bulabileceğimizi göstermeliyiz. Düşünce, sayıları
yineleyici olarak kullanıp
\[
\tuple{\num 0, F(\num 0)}, \tuple{\num 1, F(\num 1)}, \dots
\]
\emph{ikili} dizilerini hesaplamaktır. İlk ikilinin yalnızca
$\tuple{\num 0,G}$ olduğuna dikkat edin. $\tuple{\num n,F(\num n)}$ ikilisi
verildiğinde sonraki $\tuple{\num{n+1},F(\num{n+1})}$ ikilisi,
$\tuple{\num{n+1},H(\num n,F(\num n))}$ ile eşdeğer olmalıdır. Bu tek adımlı
geçişi yapan bir $T$ lambda terimi tasarlayacağız.

Ayrıntılar şöyledir. $T(u)$'yu
\[
T(u) = \tuple{S((u)_0), H((u)_0,(u)_1)}
\]
ile tanımlayın. Her $n$ sayısı için
\[
T(\tuple{\num{n}, M}) \red \tuple{\num{n+1}, H(\num{n}, M)}
\]
olduğu kolayca doğrulanır. Yukarıda önerildiği gibi $G$ ve $H$ verildiğinde
$F(u)$'yu
\[
F(u) = (u(T,\tuple{\num 0, G}))_1
\]
ile tanımlayın. Başka bir deyişle $\num n$ girdisinde $F$, $T$'yi
$\tuple{\num 0,G}$ üzerinde $n$ kez yineler, ardından ikinci bileşeni
döndürür. Başlangıçta şunlar vardır:
\begin{enumerate}
\item $\num{0} (T, \tuple{\num 0, G}) \equiv \tuple{\num{0}, G}$
\item $F(\num{0}) \equiv G$
\end{enumerate}
$n$ üzerinde tümevarımla her doğal sayı için şunları gösterebiliriz:
\begin{enumerate}
\item $\num{n+1}(T, \tuple{\num{0}, G}) \equiv \tuple{\num{n+1}, F(\num{n+1})}$
\item $F(\num{n+1}) \equiv H(\num{n}, F(\num{n}))$
\end{enumerate}
İkinci koşul için
\begin{align*}
F(\num{n+1}) & \red (\num{n+1}(T, \tuple{\num 0, G}))_1 \\
& \equiv (T(\num{n} (T, \tuple{\num 0, G})))_1 \\
& \equiv (T(\tuple{\num{n}, F(\num{n})}))_1 \\
& \equiv (\tuple{\num{n+1}, H(\num{n}, F(\num{n}))})_1 \\
& \equiv H(\num{n}, F(\num{n})).
\end{align*}
Sondan bir önceki satırda tümevarım varsayımını kullandık. İlk koşul için
\begin{align*}
\num{n+1} (T, \tuple{\num 0, G}) &
\equiv T(\num{n} (T, \tuple{\num 0, G})) \\
& \equiv T( \tuple{\num{n}, F(\num{n})}) \\
& \equiv \tuple{\num{n+1}, H(\num{n}, F(\num{n}))} \\
& \equiv \tuple{\num{n+1}, F(\num{n+1})}.
\end{align*}
Son satırda ikinci koşulu kullandık. Böylece $F(\num 0)\equiv G$ ve her $n$
için $F(\num{n+1})\equiv H(\num n,F(\num n))$ olduğunu, yani tam olarak
gerekeni göstermiş olduk.
\end{proof}

